\subsection{Stack Tecnologico}

% Ajustes locales para evitar desbordes en tablas y bloques de codigo
\setlength{\tabcolsep}{4pt}
\setlength{\LTleft}{0pt}
\setlength{\LTright}{0pt}
\lstset{basicstyle=\ttfamily\small, breaklines=true, breakatwhitespace=true, columns=fullflexible, keepspaces=true}

\subsubsection{Framework y Lenguaje}

\subsection{Flutter Framework}

\textbf{Versión:} 3.24.0 (Stable Channel) \cite{flutter_releases}

Flutter es un framework de código abierto desarrollado por Google para crear aplicaciones nativas multiplataforma desde una única base de código \cite{flutter_docs}. 

\subsubsection{Justificación de Selección}

Se seleccionó Flutter por las siguientes razones técnicas y estratégicas:

\begin{enumerate}
    \item \textbf{Desarrollo multiplataforma:} Permite compilar para Android e iOS desde un solo código base, reduciendo tiempo y costos de desarrollo en aproximadamente 40-50\% \cite{flutter_architecture}.
    
    \item \textbf{Hot Reload:} Permite ver cambios en tiempo real sin reiniciar la aplicación, acelerando significativamente el ciclo de desarrollo y debugging \cite{flutter_hot_reload}.
    
    \item \textbf{Material Design nativo:} Implementación completa de Material Design 3 sin necesidad de librerías adicionales \cite{material_design}.
    
    \item \textbf{Rendimiento nativo:} Compilación a código nativo ARM, logrando 60-120 FPS en animaciones y transiciones \cite{flutter_performance}.
    
    \item \textbf{Comunidad activa:} Más de 160,000 paquetes disponibles en pub.dev y documentación exhaustiva \cite{pub_dev}.
    
    \item \textbf{Experiencia del equipo:} Conocimiento previo en Java/Kotlin facilitó la transición a Dart.
\end{enumerate}

\subsubsection{Comparación con Alternativas}

\begin{table}[H]
\centering
\small
\begin{tabularx}{\textwidth}{|>{\raggedright\arraybackslash}p{2.8cm}|>{\centering\arraybackslash}p{2cm}|>{\centering\arraybackslash}p{2.2cm}|>{\centering\arraybackslash}p{1.8cm}|>{\centering\arraybackslash}p{1.7cm}|}
\hline
\textbf{Criterio} & \textbf{Flutter} & \textbf{React Native} & \textbf{Nativo} & \textbf{Xamarin} \\
\hline
Rendimiento & Excelente & Bueno & Excelente & Bueno \\
\hline
Hot Reload & Sí & Sí & No & Limitado \\
\hline
Curva de aprendizaje & Media & Media & Alta & Alta \\
\hline
Comunidad & Grande & Grande & Grande & Pequeña \\
\hline
UI Consistente & Sí & No & Sí & No \\
\hline
Código compartido & 100\% & 80-90\% & 0\% & 90\% \\
\hline
\textbf{Seleccion} & Si & --- & --- & --- \\
\hline
\end{tabularx}
\caption{Comparación de frameworks móviles}
\label{tab:comparacion_frameworks}
\end{table}

\textbf{Alternativas descartadas:}
\begin{itemize}
    \item \textbf{React Native:} Menor rendimiento en animaciones complejas y gráficas
    \item \textbf{Desarrollo Nativo (Kotlin/Swift):} Requiere duplicar esfuerzo para iOS
    \item \textbf{Xamarin:} Menor soporte de comunidad y documentación limitada
\end{itemize}

\subsection{Dart Language}

\textbf{Versión:} 3.5.0 \cite{dart_releases}

Dart es el lenguaje de programación oficial de Flutter, desarrollado por Google. 

\subsubsection{Características Técnicas Relevantes}

\begin{itemize}
    \item \textbf{Tipado fuerte} con inferencia de tipos automática
    \item \textbf{Null safety:} Prevención de errores NullPointerException en tiempo de compilación
    \item \textbf{Programación asíncrona:} Soporte nativo para async/await
    \item \textbf{Compilación AOT:} (Ahead of Time) para producción, genera código nativo optimizado
    \item \textbf{Compilación JIT:} (Just in Time) para desarrollo, habilita Hot Reload
    \item \textbf{Orientado a objetos:} Soporte completo para clases, interfaces y mixins
\end{itemize}

\section{Configuración del Proyecto}

\subsection{Plataforma Android}

\subsubsection{Versiones de SDK}

\begin{lstlisting}[caption=Configuración en build.gradle]
android {
    compileSdkVersion 34
    
    defaultConfig {
        applicationId "com.finara.app"
        minSdkVersion 21
        targetSdkVersion 34
        versionCode 1
        versionName "0.5.0"
    }
}
\end{lstlisting}

\begin{itemize}
    \item \textbf{minSdkVersion:} 21 (Android 5.0 Lollipop) \cite{android_versions}
    \item \textbf{targetSdkVersion:} 34 (Android 14)
    \item \textbf{compileSdkVersion:} 34 (API Level 34)
\end{itemize}

\subsubsection{Justificación de minSdkVersion 21}

\begin{table}[H]
\centering
\begin{tabularx}{\textwidth}{|l|X|}
\hline
\textbf{Razón} & \textbf{Explicación} \\
\hline
Cobertura de mercado & Cubre el 99.8\% de dispositivos Android activos según Android Distribution Dashboard \cite{android_distribution} \\
\hline
Funcionalidades modernas & Soporte completo para características de Flutter 3.x \\
\hline
Material Design & Soporte nativo para componentes de Material Design \\
\hline
Balance & Equilibrio óptimo entre compatibilidad y características modernas \\
\hline
\end{tabularx}
\caption{Justificación de minSdkVersion 21}
\end{table}

\subsubsection{Permisos Requeridos}

\begin{lstlisting}[caption=Permisos en AndroidManifest.xml]
<manifest>
    <uses-permission android:name="android.permission.INTERNET" />
    <uses-permission android:name="android.permission.CAMERA" />
    <uses-permission android:name="android.permission.READ_EXTERNAL_STORAGE" />
    <uses-permission android:name="android.permission.WRITE_EXTERNAL_STORAGE"
                     android:maxSdkVersion="28" />
</manifest>
\end{lstlisting}

\begin{description}
    \item[INTERNET:] Comunicación con backend REST API y servicios OCR externos
    \item[CAMERA:] Captura de fotografías de tickets mediante cámara del dispositivo
    \item[READ\_EXTERNAL\_STORAGE:] Acceso a galería para seleccionar imágenes existentes
    \item[WRITE\_EXTERNAL\_STORAGE:] Almacenamiento de imágenes de tickets (solo hasta Android 9)
\end{description}

\section{Dependencias Principales}

El proyecto utiliza 8 dependencias externas críticas, todas disponibles en pub.dev \cite{pub_dev}:

\small
\begin{longtable}{|p{3.9cm}|p{1.7cm}|p{7.3cm}|}
\caption{Dependencias del Proyecto FINARA} \\
\hline
\textbf{Paquete} & \textbf{Versión} & \textbf{Justificación Técnica} \\
\hline
\endfirsthead
\hline
\textbf{Paquete} & \textbf{Versión} & \textbf{Justificación Técnica} \\
\hline
\endhead
\hline
\endfoot

\texttt{fl\_chart} & 0.69.0 & Librería especializada para gráficas interactivas (barras, pastel, líneas). Seleccionada por su rendimiento (60 FPS), personalización completa y documentación clara \cite{fl_chart}. Usada en Home y Transactions. \\
\hline

\texttt{camera} & 0.10.5+9 & Plugin oficial de Flutter para acceso a cámara del dispositivo. Necesario para captura de tickets en el módulo OCR \cite{camera_plugin}. Soporta preview en tiempo real y control de flash. \\
\hline

\texttt{image\_picker} & 1.0.7 & Permite seleccionar imágenes desde galería del dispositivo. Complementa funcionalidad de cámara, permitiendo subir tickets digitales existentes \cite{image_picker}. \\
\hline

\texttt{permission\_handler} & 11.2.0 & Gestión unificada de permisos de Android/iOS en runtime. Necesario para solicitar permisos de cámara y almacenamiento de forma programática \cite{permission_handler}. \\
\hline

\texttt{shared\_preferences} & 2.2.2 & Almacenamiento local clave-valor para datos simples. Usado para autenticación (email, token), configuraciones de usuario y estado de onboarding \cite{shared_preferences}. \\
\hline

\texttt{path\_provider} & 2.1.2 & Acceso a directorios del sistema de archivos (Documents, Cache, Temp). Usado para almacenar imágenes de tickets de forma persistente \cite{path_provider}. \\
\hline

\texttt{crypto} & 3.0.3 & Funciones criptográficas (SHA-256, MD5, HMAC). Usado para hashear contraseñas antes de almacenamiento local y transmisión al backend \cite{crypto}. \\
\hline

\texttt{http} & 1.2.0 & Cliente HTTP oficial de Dart para comunicación con backend REST API. Soporta GET, POST, PUT, DELETE y manejo de headers \cite{http_package}. \\
\hline
\end{longtable}

\subsection{Justificación de Shared Preferences}

Se eligió \texttt{shared\_preferences} en lugar de SQLite local por:

\begin{table}[H]
\centering
\begin{tabularx}{\textwidth}{|p{5.2cm}|X|}
\hline
\textbf{Ventaja} & \textbf{Aplicación en FINARA} \\
\hline
Simplicidad de implementación & Prototipo no requiere base de datos compleja local \\
\hline
Velocidad de acceso & Lectura/escritura inmediata para datos de sesión \\
\hline
Suficiencia para el alcance & Almacena: email, token JWT, onboarding\_complete, setup\_data \\
\hline
Sin overhead & No requiere configuración de esquemas ni migraciones \\
\hline
Ideal para prototipos & Permite enfocarse en lógica de negocio \\
\hline
\end{tabularx}
\caption{Ventajas de SharedPreferences para FINARA}
\end{table}

\textbf{Limitaciones conocidas y planeadas:}
\begin{itemize}
    \item No escala para datos estructurados complejos o grandes volúmenes
    \item Sin soporte para queries relacionales o índices
    \item \textbf{Migración planeada:} SQLite (\texttt{sqflite}) en versión 2.0 para caché local de transacciones
\end{itemize}

\section{Archivo pubspec.yaml}

Configuración completa de dependencias:

\begin{lstlisting}[caption=pubspec.yaml del proyecto]
name: finara
description: Aplicacion de finanzas personales con OCR
version: 0.5.0+1

environment:
  sdk: ">=3.5.0 <4.0.0"

dependencies:
  flutter:
    sdk: flutter
  
  # Graficas
  fl_chart: ^0.69.0
  
  # Camara e imagenes
  camera: ^0.10.5+9
  image_picker: ^1.0.7
  
  # Permisos
  permission_handler: ^11.2.0
  
  # Almacenamiento
  shared_preferences: ^2.2.2
  path_provider: ^2.1.2
  
  # Seguridad
  crypto: ^3.0.3
  
  # Networking
  http: ^1.2.0

dev_dependencies:
  flutter_test:
    sdk: flutter
  flutter_lints: ^4.0.0

flutter:
  uses-material-design: true
\end{lstlisting}

\section{Herramientas de Desarrollo}

\begin{table}[H]
\centering
\small
\begin{tabularx}{\textwidth}{|>{\raggedright\arraybackslash}p{3.1cm}|>{\centering\arraybackslash}p{2cm}|X|}
\hline
\textbf{Herramienta} & \textbf{Versión} & \textbf{Uso} \\
\hline
Visual Studio Code & 1.85.0 & Editor principal con extensión Flutter/Dart \\
\hline
Android Studio & 2023.1 & Emuladores Android y debugging avanzado \\
\hline
Flutter DevTools & 3.24.0 & Profiling, inspector de widgets, network monitor \\
\hline
Git & 2.42.0 & Control de versiones \\
\hline
Postman & 10.20.0 & Pruebas de endpoints REST API \\
\hline
Figma & --- & Diseño de wireframes y mockups \\
\hline
\end{tabularx}
\caption{Herramientas de desarrollo utilizadas}
\end{table}
